\summarychapter{15}{Somewhere Between Heaven and Earth}{Eastern Orthodoxy}

\begin{inbrief}
Eastern Orthodoxy is the least known of the three great divisions of Christianity in the West, where
it is often taken for Roman Catholicism without the pope. The chapter explains its distinctive outlook
through the icon, its union of church and empire from Constantine to Justinian, the long conflict over
images, and the widening breach with the West that became final in 1054.
\question{What is Eastern Orthodoxy, and how did it come to differ from the Christianity of the West?}
\end{inbrief}

One summer afternoon in 1054 Cardinal Humbert and two other legates of Pope Leo IX walked into the
Church of Holy Wisdom at Constantinople as a service was about to begin, laid a bull of
excommunication on the altar and left, the cardinal shaking the dust from his feet: ``Let God look and
judge.'' A deacon ran after him begging him to take it back; it was dropped in the street. The scene
has long counted as the beginning of the great division, as Luther's theses do for Protestantism, but
it came at the end of a long process. Protestants and Catholics ask the same questions---how is a
person saved, what is the church, where is authority---and differ over the answers. In Orthodoxy the
questions themselves are different.

\begin{sidebox}{Important Events in Early Orthodoxy}
(1) Orthodoxy honors the Greek tradition of the early church: the New Testament, the Septuagint and
the first creeds are Greek, while the West was shaped by Latin. (2) It converted the Slavs; Prince
Vladimir was baptized in 988, and Orthodoxy shaped that people for a thousand years. (3) In 1054 the
Catholic and Orthodox leaders broke communion. (4) The East bore the brunt of Muslim expansion; the
crusades hurt it instead of helping; in 1453 Constantinople fell to the Turks, under whom the Orthodox
were tolerated as inferior subjects.
\end{sidebox}

\ssection{The Clue to Orthodoxy}
Orthodoxy today is about fifteen churches, mostly in Eastern Europe, joined by a common faith and
history. The best way in is not doctrine but the icon. The worshiper entering church kisses the icons
on the iconostasis, the wall of paintings between sanctuary and nave; at home icons hang in the
eastern corner and are greeted before the host. They are regarded not as human works but as windows
between earth and heaven through which the heavenly beings unite with the congregation.

The idea of image is the key. Man is made in the image of God and carries God's icon within. Western
Christians tend to see the relation of God and man in legal terms: salvation is about obtaining God's
righteousness, and the same Roman sense of law appears in the papal claim to jurisdiction. For
Orthodoxy the great theme is the incarnation and the re-creation of man. Sin does not break a legal
relationship; it wounds the divine image. Salvation is the restoration of that image by Christ, the
incarnate God: rebirth, re-creation, transfiguration. The church is not a formal institution but the
mystical body of Christ, renewed by the Spirit, a community of love in which people are made ready to
share the fellowship of Father, Son and Spirit. This is \emph{theosis}, deification. The phrase
``becoming gods'' troubles Westerners, but it does not mean becoming Father, Son or Spirit; it pictures
believers transformed into fit companions for God. As Athanasius put it, he became man so that men
might become gods.

These differences are old, but distinctively Eastern Christianity appeared under Constantine, who
first allied state and church and made purity of doctrine a concern of the empire. Some see in this
the enslavement or fall of the church; for Eastern Christians Constantine is the holy founder of the
Christian world. Rome had been turning into a religious monarchy, with the emperor as the link between
God and the world and, since Aurelian, the cult of the Invincible Sun as its religion. Then
Constantine's victory placed the emperor under the cross and in direct dependence on Christ.

\ssection{God's Will in Human Society}
Constantine was thus converted not as a man but as an emperor: through him God bound the empire to
himself. The church did not use the gospel to moderate imperial absolutism; absolutism became part of
the Christian view of the world, with the state as the bearer of religion and the expression of God's
will in society---an outlook far from the modern individual with rights that limit the state. Its
symbol was Constantinople, home of the Eastern tradition and of Byzantine civilization for a thousand
years. Constantine meant it as a holy city: in the Church of the Twelve Apostles he prepared a
thirteenth tomb for himself among the twelve, whence the title ``equal to the apostles.'' Anxious that
God would hold him responsible for Christian quarrels, he gathered the leaders of the whole church to
define the true tradition, and this became Eastern practice. From Nicaea in 325 to the seventh council
at Nicaea in 787 the emperor called and presided over every general council; the Orthodox call
themselves ``the Church of the Seven Councils.'' The councils produced the writings of the Fathers and
the creeds, but they also enlarged the role of the state. The sign of the parting ways came in 395,
when Theodosius on his deathbed divided the empire between his sons, Honorius in the West and Arcadius
in the East.

\ssection{Justinian and the Removal of Boundaries}
Under Justinian (527--565) the Byzantine blend of Roman law, Christian faith and Greek philosophy,
with a touch of the Orient, reached its height. Byzantine art subordinated the visible world to the
transcendent, nowhere more than in Hagia Sophia, the Church of Holy Wisdom, which Justinian rebuilt
and consecrated in 538, boasting that he had outdone Solomon. Its dome seemed to hang from heaven,
and its mosaics showed Constantine and Justinian offering city and church to the Mother of God.
Justinian made no distinction between Roman tradition and Christianity. The pious emperor's mission
was to keep the faith pure and protect the church, and the empire was God's instrument for guarding
and spreading Christianity. The early church had seen itself as a new people, distinct from every
natural community and belonging to the age to come. In Byzantine doctrine the state itself was the
God-given body embracing all of life, with the emperor as God's visible representative, so that the
old boundaries faded, the church dissolved into Christian society, and the sense that it is alien to
the world vanished from official thinking.

\ssection{The Conflict over Icons}
The balance of imperial power, state concern for doctrine and popular faith in icons collapsed in the
iconoclastic controversy, which lasted over a century. Beneath the dispute about pictures lay the
question of what is holy enough to be revered. From the sixth century church and state had encouraged
both icons and the veneration of holy men, not foreseeing that their multiplication would tie devotion
to local shrines and blur the line between the object and the reality it represented. There was
precedent in the emperor's image, honored as if he were present; icons of Christ and Mary gradually
replaced it, on the Bronze Gate of the palace and on coins. Leo III (717--741), perhaps reading
military defeats and an earthquake as punishment for idolatry, declared against icons after repelling
the Muslim attack on Constantinople in 717--18. The official sent to replace the Christ over the
Bronze Gate with a cross was murdered by a mob, provinces rebelled, mosaics were gouged out and a
compliant patriarch installed. The iconoclasts (``image breakers'') allowed as holy only the cross,
the Bible, the elements of the Supper, ordained clergy and dedicated buildings. Monks and the
uneducated public mostly defended the icons.

The reasoned defense came from John of Damascus, a monk in Arab-ruled Palestine and the last of the
Eastern Fathers. An image is never of the same substance as its original; it is a copy and reminder.
To deny that Christ can be pictured is to deny the incarnation. Icons are not to be worshiped but
honored like the Bible or the cross, as aids to the worship of Christ. With Patriarch Tarasius the
seventh general council, 350 bishops at Nicaea in 787, condemned iconoclasm and adopted John's view. A
synod in 843 finally confirmed it, and the Orthodox still keep the first Sunday in Lent as the Feast
of Orthodoxy.

\ssection{East Against West}
Meanwhile East and West drifted apart over an insertion in the Creed, the customs of Lent and the
bread of the Eucharist. Two ambitious men made the break permanent. In 1054 Leo IX sent the firebrand
Humbert with terms sure to fail; the patriarch Michael Cerularius answered his intolerance in kind,
and Humbert left the bull on the altar. Later the empire of Constantine and Justinian was whittled
away by barbarians and Islam until Constantinople itself fell to the Turks in 1453. Christians became
a minority under Muslim rule and looked to the patriarch as their political head. Today only Greece
and half of Cyprus remain Greek Orthodox. But a corridor lay open to the north: Boris of Bulgaria was
converted in the ninth century and Vladimir of Kiev in the tenth. His envoys, after the liturgy in
Hagia Sophia, reported, ``We know not whether we were in heaven or on earth.'' Russia made Orthodoxy's
beauty its own, and Moscow came to see itself as the ``third Rome'' after Rome and Constantinople, its
ruler bearing Caesar's title, \emph{tsar}.

\begin{keyfacts}
\keypeople{Cardinal Humbert; Pope Leo IX; Michael Cerularius; Constantine; Theodosius; Justinian; Leo
III; John of Damascus; Patriarch Tarasius; Boris of Bulgaria; Vladimir of Kiev; Athanasius.}
\keydates{%
  \kd{395}{Theodosius divides the empire between East and West}
  \kd{527--565}{Reign of Justinian}
  \kd{538}{Hagia Sophia consecrated}
  \kd{717--741}{Reign of the emperor Leo III; the attack on icons begins}
  \kd{787}{Seventh general council (Nicaea) approves the veneration of icons}
  \kd{843}{Final restoration of the icons: the Feast of Orthodoxy}
  \kd{988}{Baptism of Vladimir of Kiev}
  \kd{1054}{Humbert's bull of excommunication: break between East and West}
  \kd{1453}{Constantinople falls to the Turks}}
\keyterms{Icon; iconostasis; image of God; \emph{theosis} (deification); ``equal to the apostles'';
Church of the Seven Councils; Byzantine; iconoclasm; Feast of Orthodoxy; third Rome.}
\end{keyfacts}
